\documentclass[11pt]{article}

\usepackage[margin=1.1in]{geometry}
\usepackage[T1]{fontenc}
\usepackage[utf8]{inputenc}
\usepackage{lmodern}
\usepackage{microtype}
\usepackage{setspace}
\usepackage[colorlinks=true, linkcolor=black, urlcolor=blue]{hyperref}

\setstretch{1.15}
\setlength{\parskip}{0.4em}

\title{\textbf{Waves, Particles, and Sound}\\
\large An Essay Companion to the Virtual Exhibition}
\author{Yuval Shemla\\ \small Columbia University --- Music Humanities --- Summer A 2026}
\date{July 2026}

\begin{document}

\maketitle

\begin{abstract}
\noindent This essay accompanies \emph{Waves, Particles, and Sound}, a browser-based
three-dimensional exhibition examining ten episodes, from 1863 to 2023, in which developments
in physics altered the materials, methods, or meaning of Western music. Ten works stand in a
circular gallery in chronological order, each paired with the physics that conditioned it.
For every station, this essay presents the physics, the music, and the curatorial reasoning:
why the work was included, what it represents, and how it relates to the sequence as a whole.
\end{abstract}

\newpage
\tableofcontents
\newpage

%=====================================================================
\section{Introduction}
%=====================================================================

This exhibition examines ten episodes in which a development in physics --- in acoustics,
electromagnetism, statistical mechanics, nuclear physics, or quantum theory --- preceded and
conditioned a change in Western compositional practice. The earliest episode begins with
Hermann von Helmholtz's acoustical measurements of 1863; the latest ends with a duet between
a violinist and a quantum computer in 2023.

The exhibition takes the form of a rotunda. The visitor begins at the center of a circular
hall, surrounded by a ring of ten dates --- 1894 to 2023 --- displayed above ten station
walls. Between the first date and the last stands a single door, at which the sequence both
begins and ends. The circular plan is not decorative: the final stations return, with new
analytical tools, to questions raised by the first ones, so that the geometry of the room
reproduces the structure of the argument.

A methodological note is required. The connections between physics and music presented here
are not all of one kind, and each station identifies its own. Some are \emph{direct}: a
composer works from the equations themselves (Xenakis, Grisey, Miranda). Some are
\emph{instrumental}: the physics constitutes the instrument (Theremin, Lucier). Some are
\emph{experiential} or \emph{cultural}: a discovery alters the prevailing description of
nature, and artists respond in parallel without documented contact (Ives, Stravinsky). One is
a \emph{structural analogy} argued retrospectively by a physicist (Cage), and one is a
\emph{moral response} to what physics made possible (Penderecki). An exhibition that
presented every pairing as direct influence would overstate its evidence. Distinguishing the
categories is part of the argument.

%=====================================================================
\section{Station 1 --- The Anatomy of a Single Tone (1894)}
%=====================================================================

\subsection{The Physics}
In 1863, the German physicist Hermann von Helmholtz published \emph{On the Sensations of
Tone}, proving that every musical tone is a composite: a fundamental pitch plus a ladder of
overtones. The recipe of overtones gives each instrument its unique color, its
\emph{timbre}. Using tuning forks and glass resonators, Helmholtz revealed the hidden
architecture inside every note.

\subsection{The Music}
In 1894, Claude Debussy premiered \emph{Pr\'elude \`a l'apr\`es-midi d'un faune}. The piece
opens with a sinuous flute solo floating free of any clear key. Debussy layers translucent
washes that shimmer and dissolve rather than march through conventional harmony. He treats
chords as pure \emph{sonorities} --- blocks of color valued for how they sound, not where
they lead.

\subsection{Why This Work Is in the Exhibition}
Helmholtz showed that timbre arises from the overtone series. Debussy, who studied acoustics
at the Paris Conservatoire and knew Helmholtz's work, composed with that structure in mind
--- hearing chords as textures rather than functions. The connection is \emph{informed
influence}.

Boulez dated the beginning of modern music to the opening flute solo of the \emph{Faun}. The
exhibition opens here because the work is the clearest documented case of scientific
knowledge altering compositional practice without appearing explicitly in the score: Debussy
encountered acoustics during his Conservatoire training and subsequently treated sonority ---
color, texture, resonance --- as primary material rather than as harmonic grammar. The
station also establishes the exhibition's structural through-line. The overtone series
Helmholtz measured recurs twice in the circle: as the resonant frequencies of Lucier's room
(Station 8) and as the spectrum Grisey orchestrates (Station 9). A single acoustical
discovery thus spans three generations of compositional technique.

%=====================================================================
\section{Station 2 --- Sound in Motion (1908)}
%=====================================================================

\subsection{The Physics}
In 1842, Christian Doppler proposed that wave frequency changes when source and observer move
relative to each other. An ambulance siren rises as it approaches and drops as it recedes ---
the waves compress in front of the moving source and stretch behind it. The effect was first
tested with trumpeters on a moving train in 1845.

\subsection{The Music}
Charles Ives grew up with a father who positioned marching bands on opposite sides of town,
each playing different music. From this childhood he composed \emph{The Unanswered Question}
(1908): three groups placed apart in space --- offstage strings, a solo trumpet posing ``the
perennial question,'' a woodwind quartet answering with increasing frenzy --- each in its own
musical world, overlapping but never converging.

\subsection{Why This Work Is in the Exhibition}
Ives's father's experiments were Doppler's principle in action: sound changes with spatial
motion. Ives composed from the listener's moving body, creating music where spatial position
shapes perception. The link is \emph{experiential}.

The pairing originates in this course: the relation between the Doppler effect and Ives's
spatial practice was drawn in lecture, and the station develops that observation. Beyond the
course connection, Ives is the first composer in the sequence to treat spatial position as a
compositional parameter. The offstage strings of \emph{The Unanswered Question} anticipate a
century of spatialized music, from antiphonal seating arrangements to electroacoustic
diffusion --- and, not incidentally, the navigable virtual space of the exhibition itself.
The work's program, a question posed and never satisfactorily answered, also states in
miniature the question the exhibition examines: what kind of answer music gives to a changed
description of nature.

%=====================================================================
\section{Station 3 --- Simultaneous Perspectives (1913)}
%=====================================================================

\subsection{The Physics}
In 1905, Einstein published Special Relativity: there is no single, absolute frame of
reference. Time itself is relative; simultaneity is an illusion. In the same decade, Picasso
shattered the single viewpoint in painting --- his 1912 \emph{Guitar} shows an object from
multiple angles at once. Reality depended on where you stood.

\subsection{The Music}
On May 29, 1913, Stravinsky's \emph{The Rite of Spring} caused a riot in Paris. He stacked
rhythms that refuse to align (\emph{polyrhythm}), layered clashing keys
(\emph{polytonality}), and drove pounding asymmetric accents. Multiple musical ``realities''
coexist simultaneously, each with its own pulse and tonal center.

\subsection{Why This Work Is in the Exhibition}
Einstein, Picasso, and Stravinsky never collaborated. But within one decade, physics, visual
art, and music each dismantled the idea of a single, stable perspective. Relativity denied
absolute time; Cubism denied a single viewpoint; \emph{The Rite} denied a single meter and
key. The connection is \emph{synchronous rupture}.

Not every pairing in the exhibition claims direct influence, and this station defines the
alternative. There is no documented link between Einstein and Stravinsky; Picasso's
\emph{Guitar} occupies the center of the wall as a contemporaneous parallel, not as evidence
of exchange. The station instead documents synchrony: within roughly one decade, physics,
painting, and music each abandoned the premise of a single privileged frame of reference. Its
inclusion serves a methodological purpose --- it marks the boundary between demonstrable
influence and cultural simultaneity, and keeps the exhibition's stronger claims credible. The
premiere's notoriety indicates how audible such a conceptual shift was to its first
audiences.

%=====================================================================
\section{Station 4 --- Music from Thin Air (1920)}
%=====================================================================

\subsection{The Physics}
In 1920, the Russian physicist Leon Theremin was researching high-frequency circuits when he
noticed that moving his hand near an antenna changed the audible pitch. His body acted as a
variable capacitor --- \emph{heterodyne oscillation} produced an audible difference tone from
two radio-frequency signals. He had stumbled upon the first instrument you play without
touching.

\subsection{The Music}
Clara Rockmore, a Lithuanian-born violin prodigy, became the theremin's undisputed virtuoso.
She coaxed from it a tone part cello, part soprano voice --- a continuous, vibrato-rich
legato unlike anything else in music. Her right hand sculpted pitch in the air; her left
controlled volume.

\subsection{Why This Work Is in the Exhibition}
The theremin does not \emph{use} electromagnetic principles as metaphor --- it \emph{is} an
electromagnetic instrument. The performer's body becomes a variable capacitor in a circuit.
Leon Theremin was a physicist first and an instrument builder second. The connection is
\emph{direct and instrumental}.

The station has a biographical origin: an encounter with a theremin exhibited at the
Metropolitan Museum of Art, which suggested the premise of this exhibition. Its justification
is broader. The theremin is the first performable instrument without a vibrating body; its
sound is produced entirely by the interaction of a performer with an electromagnetic field,
making it the most literal case in the exhibition of physics functioning as an instrument.
Historically it stands at the head of the electronic lineage that runs through the
synthesizer to the present. Rockmore's performance practice is included to document that the
instrument sustained serious interpretive standards from its first generation.

%=====================================================================
\section{Station 5 --- The Sound of Nothing (1952)}
%=====================================================================

\subsection{The Physics}
Quantum mechanics revealed that perfect emptiness does not exist. Even in a vacuum, virtual
particles flicker into existence and annihilate on timescales too brief to observe, but their
effects are measurable --- the \emph{Casimir effect} pushes two plates together using nothing
but vacuum fluctuations. ``Nothing'' is not nothing.

\subsection{The Music}
In 1951, John Cage visited Harvard's anechoic chamber expecting silence. Instead he heard two
sounds: his nervous system and his blood circulating. Silence was unreachable. In 1952 he
premiered \emph{4'33''}: three movements during which the performer plays nothing. The
``music'' is whatever ambient sounds occur --- and every performance is unique. In 2020 the
Berlin Philharmonic performed it to an empty, locked-down hall.

\subsection{Why This Work Is in the Exhibition}
The physicist Donald Spector, in a paper titled ``John Cage Adores a Vacuum,'' argued that
\emph{4'33''} is the auditory analogue of the quantum vacuum: in both cases, the absence of
deliberate signal reveals that ``empty'' space is full of activity. The connection is
\emph{structural analogy}.

This station occupies the conceptual center of the exhibition. The question it poses ---
whether silence exists --- receives parallel negative answers from quantum field theory and
from Cage's experience in the anechoic chamber, and the parallel has been argued formally in
Spector's published paper. The pairing is an analogy rather than an influence, and is labeled
as such. One curatorial decision follows from the work itself: the station provides no
recording, since the composition consists of the ambient sound present at each performance.
Cage's redefinition of the musical work --- from organized sound to framed attention ---
conditions every station that follows.

%=====================================================================
\section{Station 6 --- Music by Probability (1955)}
%=====================================================================

\subsection{The Physics}
In the 1860s, Maxwell and Boltzmann developed statistical mechanics --- describing billions
of gas molecules not individually but through probability distributions. Brownian motion, the
random jiggling of pollen grains buffeted by invisible molecules, became a visible signature
of this statistical reality. Order emerges from randomness.

\subsection{The Music}
Iannis Xenakis --- trained as an engineer and architect under Le Corbusier --- applied
statistical mechanics directly in \emph{Pithoprakta} (1955--56). Each of 46 string players
performs an independent glissando with pitch, speed, and direction determined by
Maxwell--Boltzmann equations. The result is a seething, swarming mass --- the sonic
equivalent of a gas cloud.

\subsection{Why This Work Is in the Exhibition}
This is \emph{explicit mathematical application}. Xenakis imported the equations into his
compositional method and published them in \emph{Formalized Music} (1971). The orchestra
becomes a gas: 46 independent particles whose collective motion produces emergent form.

\emph{Pithoprakta} is the exhibition's strongest case of direct mathematical transfer.
Xenakis did not analogize statistical mechanics; he calculated with the Maxwell--Boltzmann
distribution and documented the method in \emph{Formalized Music}. The station represents the
professional convergence the exhibition traces --- Xenakis worked simultaneously as engineer,
architect, and composer --- and it establishes the precondition for the final stations: once
compositional decisions can be generated from mathematical models, computation becomes a
musical instrument. The title translates as ``actions by means of probability.'' Subsequent
algorithmic and stochastic composition, including the quantum-generated music of Station 10,
develops from the procedure demonstrated here.

%=====================================================================
\section{Station 7 --- Sound After the Unthinkable (1960)}
%=====================================================================

\subsection{The Physics}
On July 16, 1945, physicists detonated the first nuclear weapon in the New Mexico desert.
Nuclear fission --- splitting atomic nuclei --- released energy equivalent to 21,000 tons of
TNT. Weeks later, bombs fell on Hiroshima and Nagasaki, killing over 200,000 people.

\subsection{The Music}
Penderecki's \emph{Threnody for the Victims of Hiroshima} (1960) uses no melody and no
harmony. He invented a graphic notation --- squiggles, blacked-out blocks, time-space
diagrams --- to direct 52 strings in \emph{tone clusters}, col legno strikes, and unearthly
quarter-tone glissandi: a wall of anguished sound that tightens, shrieks, and fades to
silence.

\subsection{Why This Work Is in the Exhibition}
The connection is \emph{moral response}. Penderecki used music to bear witness to what
nuclear physics had made possible.

The other stations document what physics contributed to music; this one documents music's
response to what physics made possible. The \emph{Threnody} is included as the exhibition's
necessary counterweight: an account of the century's physics that omitted its destructive
consequences would be incomplete. The connection is one of moral response rather than
technique, though a technical parallel exists --- the graphic notation Penderecki devised for
the work's clusters and glissandi visually resembles particle-track diagrams, an accidental
convergence between the notation of catastrophe and the notation of matter. Within the
circle, this station marks the point at which the exhibition's subject acquires ethical
weight.

%=====================================================================
\section{Station 8 --- The Room Reveals Itself (1969)}
%=====================================================================

\subsection{The Physics}
Every enclosed space has resonant frequencies --- pitches at which sound waves reflect off
walls in patterns that reinforce themselves, creating standing waves. The room acts as a
filter, strengthening some frequencies and suppressing others. Every room has a sonic
fingerprint.

\subsection{The Music}
In 1969, Alvin Lucier recorded himself speaking: ``I am sitting in a room different from the
one you are in now\ldots'' He played the tape back into the room, re-recorded it, and
repeated the cycle 32 times. Consonants blurred, vowels smeared, and his voice --- including
his stutter --- dissolved into pure shimmering tones: the resonant frequencies of the room
itself.

\subsection{Why This Work Is in the Exhibition}
The composition \emph{is} the physics. There is no metaphor, no translation. Each
re-recording cycle strips away frequencies the room does not reinforce and amplifies those it
does. The room becomes the instrument, the score, and the performer. The connection is
\emph{identity}.

The work is included because it is not a representation of an acoustical process but an
execution of one. The score is a set of procedural instructions; the apparatus is two tape
recorders; the result is determined by the resonant frequencies of the performance space. The
piece is also uniquely reproducible --- it can be realized in any room, with different
results in each --- which makes it the exhibition's clearest demonstration that the acoustics
of Station 1 apply beyond instruments. The overtone series returns here as a property of
architecture, completing one of the circle's internal threads.

%=====================================================================
\section{Station 9 --- Orchestrating the Spectrogram (1975)}
%=====================================================================

\subsection{The Physics}
In 1822, Joseph Fourier proved that any complex waveform can be decomposed into simple sine
waves of different frequencies and amplitudes. Applied to sound, this means any tone --- a
trombone blast, a vowel, a cymbal crash --- can be displayed as a \emph{spectrogram}: a
visual map of its frequency components.

\subsection{The Music}
In 1975, G\'erard Grisey recorded a low E on a trombone and studied its spectrogram at IRCAM.
Then he \emph{orchestrated the spectrogram}: each overtone assigned to a different instrument
--- fundamental to trombone, second partial to horn, third to clarinet, upward through the
ensemble. \emph{Partiels} begins as a single trombone note and blooms into a shimmering chord
that is literally the inside of that note made audible.

\subsection{Why This Work Is in the Exhibition}
This is the most direct physics-to-music translation in the exhibition. Grisey did not
interpret Fourier analysis --- he fed sound into a machine, read the data, and transcribed it
onto orchestral staves. The ensemble \emph{plays} a spectrogram. The school this founded is
called \emph{Spectral Music}. The connection is \emph{direct transcription}.

\emph{Partiels} closes the circle the exhibition opens at Station 1: it returns to the
interior structure of a single tone, now analyzed by Fourier methods at IRCAM rather than
measured with resonators. The station documents the most direct transcription in the
exhibition --- spectral data transferred to orchestral staves --- and its consequences: the
spectral school that \emph{Partiels} inaugurated became one of the most widely taught
compositional movements of the late twentieth century. The rotunda's circular plan reflects
this relation: the end of the sequence returns to the material of its beginning.

%=====================================================================
\section{Station 10 --- Quantum Duet (2023)}
%=====================================================================

\subsection{The Physics}
A quantum computer uses \emph{qubits} --- quantum bits that exist in superposition of both 0
and 1 simultaneously. Only when you measure a qubit does it collapse to a definite value, and
the outcome is fundamentally probabilistic. In 2022, IBM unveiled the Eagle processor with
127 qubits --- powerful enough for its outputs to drive creative decisions.

\subsection{The Music}
In Eduardo Reck Miranda's \emph{Qubism} (2023), a violinist plays a phrase; an IBM quantum
computer processes it through a quantum circuit, and the probabilistic output of qubit
measurements determines the responding phrase. No two performances are identical. The form of
call-and-response --- heard in this course from spirituals to jazz --- meets the newest
technology available to music.

\subsection{Why This Work Is in the Exhibition}
The 160-year arc closes here. Helmholtz used tuning forks; Miranda uses qubits. In
\emph{Qubism}, quantum computation \emph{generates} musical decisions through measurement and
superposition. The connection is \emph{generative and direct}.

The sequence ends at the present. \emph{Qubism} is included because it inverts the relation
documented in the preceding stations: physics appears here not as subject matter, method, or
instrument, but as a performing partner whose responses are generated by quantum measurement.
The work also continues a lineage examined throughout this course --- call and response, from
the spirituals onward --- in its newest technological form. Its indeterminacy is structural:
no two performances are identical, because the underlying measurements are probabilistic. The
exhibition therefore closes without resolution, and the circle returns the visitor to the
door at which it began.

%=====================================================================
\section{Epilogue}
%=====================================================================

The sequence ends where it began. Across the ten stations a pattern recurs: a development in
physics --- in acoustics, electromagnetism, statistical mechanics, or quantum theory ---
precedes and conditions a change in compositional practice.

The pattern has not closed. Machine-learning systems now generate music; gravitational-wave
observatories render astrophysical events as audible signals. The material for an eleventh
station already exists. What that station will hold is the question the exhibition leaves
open.

%=====================================================================
\section*{A Note on the Form}
\addcontentsline{toc}{section}{A Note on the Form}
%=====================================================================

The exhibition itself is a browser-based three-dimensional rotunda built with A-Frame,
navigable with keyboard and mouse. Ten station walls circle the visitor in chronological
order beneath a ring of dates; a single monumental door marks the beginning and end of the
sequence. Each station presents concise wall text, framed images of the physics and the
music, and an audio excerpt of the work that begins automatically as the visitor arrives; a
``Read More'' letter carries the fuller discussion summarized in this essay. The exhibition
runs locally from a single self-contained HTML file.

\end{document}
